\documentclass[aspectratio=169,11pt]{beamer}
\usetheme{default}
\setbeamertemplate{navigation symbols}{}
\usepackage{graphicx}
\usepackage[T1]{fontenc}
\definecolor{ncsured}{HTML}{CC0000}
\definecolor{ncsudk}{HTML}{990000}
\setbeamercolor{structure}{fg=ncsured}
\setbeamercolor{frametitle}{fg=white,bg=ncsured}
\setbeamercolor{title}{fg=ncsudk}
\setbeamercolor{itemize item}{fg=ncsured}
\setbeamercolor{itemize subitem}{fg=ncsured}
\setbeamerfont{frametitle}{series=\bfseries}
\setbeamertemplate{frametitle}{%
  \nointerlineskip\vskip2pt%
  \begin{beamercolorbox}[wd=\paperwidth,ht=2.4ex,dp=1ex,leftskip=0.3cm]{frametitle}%
  \bfseries\insertframetitle%
  \end{beamercolorbox}}
\setbeamertemplate{footline}{%
  \hbox{\begin{beamercolorbox}[wd=\paperwidth,ht=2.2ex,dp=1ex,leftskip=.3cm,rightskip=.3cm]{}%
  \tiny Optimal Market Making for Cryptocurrency Options \hfill Week 3 \hfill \insertframenumber%
  \end{beamercolorbox}}}

\title{Week 3: The Avellaneda-Stoikov Model}
\subtitle{A worked example of optimal market making via stochastic control}
\date{North Carolina State University \,\textbar\, Summer 2026}
\titlegraphic{\includegraphics[width=4.2cm]{/Users/dendisuhubdy/Github/ncsu/assets/ncsu_logo.png}}
\begin{document}
\begin{frame}[plain]\titlepage
\begin{center}\tiny Industry Mentor: Dendi Suhubdy (Bitwyre)  ·  Guest: Fenni Kang (Coincall)\end{center}\end{frame}
\begin{frame}{Learning Objectives}
\begin{itemize}
  \item Set up the Avellaneda-Stoikov control problem with CARA utility
  \item Understand the HJB equation and its reduction via the ansatz
  \item Interpret the reservation price and optimal half-spread
\end{itemize}\end{frame}
\begin{frame}{Lecture 1: Setup and Dynamics}
\begin{itemize}
  \item Mid-price as arithmetic Brownian motion (chosen for tractability)
  \item Poisson order flow with intensity lambda(d) = A exp(-kappa d)
  \item Wealth and inventory dynamics under continuous quoting
  \item Exponential (CARA) utility decouples value from wealth
\end{itemize}\end{frame}
\begin{frame}{Lecture 2: HJB and Closed-Form Quotes}
\begin{itemize}
  \item Dynamic programming principle -\textgreater{} HJB equation
  \item Separation ansatz V = -exp(-gamma(x+qS)) theta(t,q)
  \item Reservation price r = S - q gamma sigma\textasciicircum{}2 (T-t)
  \item Half-spread = (1/gamma) ln(1+gamma/kappa) + (gamma sigma\textasciicircum{}2 / 2)(T-t)
\end{itemize}\end{frame}
\begin{frame}{Recitation}
\begin{itemize}
  \item Implement the closed-form quotes in PyTorch
  \item Visualize reservation price and spread vs. inventory
  \item Run a risk-aversion (gamma) sensitivity sweep
\end{itemize}\end{frame}
\begin{frame}{Reading \& Problem Set 3 (Python/PyTorch)}
\begin{itemize}
  \item Avellaneda-Stoikov (2008); Cartea-Jaikumar-Penalva Ch. 10
  \item PS3: ODE solver, closed-form quotes, simulation, gamma sweep
\end{itemize}\end{frame}
\end{document}